\section{Process learning in the structured Transformer}
\label{app:joint}
\label{app:step-supplied-routing}



\subsection{Model and exact features}

We count modulo $K$ with a fixed query of probability $p\in(0,1)$ and a uniform start; every non-query letter has the same embedding. Let $\bm x_t=\bm e_{s_{t-1}}$, $\mathbf R\bm e_s=\bm e_{s+1}$, $r_b=\sigma\bigl(-(-1)^ba\bigr)$ for $b\in\{0,1\}$ and $\sigma(z)=(1+e^{-z})^{-1}$.

\paragraph{Seed.} Fix orthonormal residual directions $\hat{\bm e}_0,\hat{\bm e}_\ell$, a count register $(\hat{\bm e}^{\mathrm c}_s)_{s\in\mathbb Z_K}$, a gate register
$(\hat{\bm e}^{\mathrm g}_s)_{s\in\mathbb Z_K}$, and a unit vector $\hat{\bm e}_{\mathrm h}$ of the head space. The embeddings are
\[
 E(s)=\initscale\,\hat{\bm e}^{\mathrm c}_s,\qquad E(\text{query letter})=\initscale(\hat{\bm e}_0-\hat{\bm e}_\ell),\qquad E(\text{other letters})=\initscale(\hat{\bm e}_0+\hat{\bm e}_\ell).
\]
The gate head is head $1$ of block $2$. At the position $i_t$ that holds $s_{t-1}$ and predicts $s_t$ ($t\ge2$), its fixed mask admits
$i_t$ and the position $j_t$ of the current letter, and its relative bias is zero. Its maps are $\mathbf W_Q=q\,\bm\Pi_Q$, $\mathbf W_K=k\,\bm\Pi_K$ and
$\mathbf W_V=v\,\bm\Pi_V$, with scalars $q,k,v$ and fixed matrices
\[
 \bm\Pi_Q=\textstyle\sum_s\hat{\bm e}_{\mathrm h}(\hat{\bm e}^{\mathrm c}_s)^{\!\top},\qquad \bm\Pi_K=\hat{\bm e}_{\mathrm h}\hat{\bm e}_\ell^{\top},\qquad \bm\Pi_V=\sum_s\hat{\bm e}^{\mathrm g}_s(\hat{\bm e}^{\mathrm c}_s)^{\!\top}.
\]
All other embeddings, heads and MLP weights are zero, and the MLP biases are $\pm\initscale$. The only trainable readout blocks are the
$K\times K$ blocks $\mathbf U$ (count register) and $\mathbf W=w(\mathbf I-\mathbf R)$ (gate register), both initialized at zero.
The query and key train along $\bm\Pi_Q,\bm\Pi_K$; the value map $v\bm\Pi_V$ is fixed (\cref{rem:J-line}).

\paragraph{Training.} The trained parameters are $\bm\vartheta=(\mathbf U,\mathbf W,q,k)$, with $\mathbf W$ on the line $\R(\mathbf I-\mathbf R)$, and
$\ddt{\bm\vartheta}=-\nabla_{\bm\vartheta}\Lstep(\bm\vartheta)$ for the Euclidean (Frobenius) norm on $\R^{K\times K}\times\R(\mathbf I-\mathbf R)\times\R\times\R$. The value factor $v$
and all other parameters are fixed. The four trained parameters follow this one flow at all times; there are no stages and no stopping
rule. The initialization is
 $\mathbf U(0)=0,\qquad w(0)=0,\qquad \tfrac{1}{2}\initscale\le|q(0)|,\ |k(0)|,\ |v|\le2\initscale,$ 
with arbitrary signs.

\begin{lemma}[Exact features]
\label{lem:J-features}
For every $\bm\vartheta$ and $t\ge2$, the logits of the prediction of $s_t$ are
 $\bm z_t=\initscale \mathbf U\bm x_t+\initscale v\,r_{b_t}\,\mathbf W\bm x_t,\qquad a=\frac{\initscale^2qk}{\sqrt{d_{\mathrm h}}}.$ 
\end{lemma}

\begin{proof}
Block $1$ and all MLPs add zero, so each residual is its embedding. The query at $i_t$, the keys and the values are
\begin{gather*}
 \text{query}_{i_t}=\initscale q\,\hat{\bm e}_{\mathrm h},\qquad \text{key}_{j_t}=(-1)^{b_t}\initscale k\,\hat{\bm e}_{\mathrm h},\qquad \text{key}_{i_t}=0,\\
 \text{value}_{i_t}=\initscale v\,\hat{\bm e}^{\mathrm g}_{s_{t-1}},\qquad \text{value}_{j_t}=0 .
\end{gather*}
The scores are $(-1)^{b_t}a$ at $j_t$ and $0$ at $i_t$, so $i_t$ gets weight $r_{b_t}$. The residual at $i_t$ becomes
 $\initscale\hat{\bm e}^{\mathrm c}_{s_{t-1}}+\initscale v\,r_{b_t}\,\hat{\bm e}^{\mathrm g}_{s_{t-1}},$ 
and the readout gives $\bm z_t$.
\end{proof}

\begin{lemma}[Cyclic symmetry]
\label{lem:J-symmetry}
$\mathbf R\mathbf U(t)=\mathbf U(t)\mathbf R$ for all $t\ge0$; $\mathbf W(t)\in\R(\mathbf I-\mathbf R)$ commutes with $\mathbf R$ by definition.
\end{lemma}

\begin{proof}
Shifting $s_0$ by one preserves the data law and the match bits, and $\mathrm{CE}(\bm z,j+1)=\mathrm{CE}(\mathbf R^{-1}\bm z,j)$ for the cross-entropy. By \cref{lem:J-features},
 $\Lstep\circ S=\Lstep,\qquad S(\mathbf U,\mathbf W,q,k)=(\mathbf R^{-1}\mathbf U\mathbf R,\ \mathbf W,\ q,\ k).$ 
$S$ is a linear isometry, so $S\bm\vartheta(t)$ solves the flow, with $S\bm\vartheta(0)=\bm\vartheta(0)$. By uniqueness, $S\bm\vartheta(t)=\bm\vartheta(t)$.
\end{proof}

\subsection{Reduction}

Let $N=K-2$, $\rho=\frac{2}{K}$, $\pi_0=1-p$, $\pi_1=p$, $H_p=-p\log p-(1-p)\log(1-p)$ and $p_{\min}=\min(p,1-p)$. For $m,s\in\R$ let
 $F(m;s)=\log\bigl(1+e^{-m}+N e^{-s-m/2}\bigr),\qquad \phi(m;s)=-\pd{F}{m}(m;s);$ 
for $K=2$, $F(m;s)=\ell(m):=\log(1+e^{-m})$.

\begin{lemma}[Reduced loss and flow]
\label{lem:J-reduce}
Write $\mathbf U=\sum_ju_j\mathbf R^j$. Then $u_j=u_{\mathrm{bg}}$ for $2\le j\le K-1$. Let
 $c=\initscale(u_0-u_1),\qquad s=\initscale\bigl(\tfrac{1}{2}(u_0+u_1)-u_{\mathrm{bg}}\bigr),\qquad g=2\initscale vw,\qquad m_b=(-1)^b\bigl(c+g\,r_b\bigr),$ 
with $s=0$ if $K=2$. The prediction of $s_t$ with match bit $b_t=b$ puts logits $s+\frac{m_b}{2}$ on its target, $s-\frac{m_b}{2}$ on the other count of
$\{s_{t-1},s_{t-1}+1\}$, and $0$ on the remaining $N$ counts, up to a common shift. Hence, for every $n$,
 $\Lstep=\mathcal L(c,g,s,a):=(1-p)\,F(m_0;s)+p\,F(m_1;s).$ 
With $\qkbalance=q(0)^2-k(0)^2$,
\begin{equation}
 \begin{gathered}
 \ddt{c}=-\rho\initscale^2\pd{\mathcal L}{c},\qquad \ddt{g}=-\rho\initscale^2v^2\pd{\mathcal L}{g},\qquad \ddt{s}=-\frac{\initscale^2}{2N}\pd{\mathcal L}{s}\ (K\ge3),\qquad \ddt{a}=-\frac{\initscale^4S}{d_{\mathrm h}}\pd{\mathcal L}{a},\\
 S:=q^2+k^2=\sqrt{\qkbalance^2+\frac{4d_{\mathrm h}a^2}{\initscale^4}}.
 \end{gathered}
 \label{eq:J-flow}
\end{equation}
\end{lemma}

\begin{proof}
\emph{Background coordinates.} By \cref{lem:J-features}, $\Lstep$ depends on $\mathbf U$ through $(u_j)$, and the targets are at offsets $0$ and $1$.
Permuting $u_2,\ldots,u_{K-1}$ preserves $\Lstep$ and $\|\mathbf U\|_F^2=K\sum_ju_j^2$ and fixes $\mathbf U(0)=0$; as in \cref{lem:J-symmetry},
$u_j=u_{\mathrm{bg}}$ for $j\ge2$.

\emph{Logits.} For $\bm x_t=\bm e_{s'}$, $\mathbf W\bm x_t=w(\bm e_{s'}-\bm e_{s'+1})$. By \cref{lem:J-features}, after subtracting $\initscale u_{\mathrm{bg}}$,
 $z_{t,s'}=s+\tfrac{1}{2}(c+g\,r_{b_t}),\qquad z_{t,s'+1}=s-\tfrac{1}{2}(c+g\,r_{b_t}),\qquad z_{t,s'+j}=0\ \ (2\le j\le K-1).$ 
The target is $s'+b_t$, so its logit is $s+\frac{m_{b_t}}{2}$ and the other one is $s-\frac{m_{b_t}}{2}$. The cross-entropy is
$\log(1+e^{-m}+Ne^{-s-m/2})=F(m;s)$.

\emph{Flow.} At a circulant $\mathbf U$ the Frobenius gradient $\nabla_{\mathbf U}\Lstep$ is circulant (it is invariant under the isometry $S$ of
\cref{lem:J-symmetry}), so the flow stays on circulant matrices and acts on their coefficients as $\ddt{u_j}=-\frac{1}{K}\pd{\Lstep}{u_j}$ and $\ddt{w}=-\frac{1}{2K}\pd{\Lstep}{w}$, since
$\|\mathbf U\|_F^2=K\sum_ju_j^2$ and $\|\mathbf W\|_F^2=2Kw^2$. Substituting the definitions of $c,g,s$ in \cref{lem:J-reduce}
into these metric gradients gives their equations in \cref{eq:J-flow}.
For the gate,
\begin{gather*}
 \ddt{q}=-\frac{\initscale^2k}{\sqrt{d_{\mathrm h}}}\pd{\mathcal L}{a},\qquad \ddt{k}=-\frac{\initscale^2q}{\sqrt{d_{\mathrm h}}}\pd{\mathcal L}{a},\qquad \frac{d}{dt}(q^2-k^2)=0.
\end{gather*}
Differentiating $a=\initscale^2qk/\sqrt{d_{\mathrm h}}$ gives the remaining equation of \cref{eq:J-flow}; $S^2=(q^2-k^2)^2+4q^2k^2$ gives its expression for $S$. The flow exists for all $t\ge0$: the vector field is smooth, $|\pd{\mathcal L}{c}|,|\pd{\mathcal L}{g}|,|\pd{\mathcal L}{s}|\le1$, so
$c,g,s$ grow at most linearly, and $|\pd{\mathcal L}{a}|\le\frac{|g|}{4}$, so $(q,k)$ grows at most exponentially by Gr\"onwall.
\end{proof}

\paragraph{Notation.} $r_0+r_1=1$ and $\tau:=r_1-r_0=\tanh(\frac{a}{2})$. Let
\begin{gather*}
 \mathrm{err}_b=\pi_b\,\phi(m_b;s),\qquad \mathrm{err}=\mathrm{err}_0+\mathrm{err}_1,\\
 h=\tfrac{1}{2}(m_0-m_1)=c+\tfrac{g}{2},\qquad R=m_0+m_1=-g\tau,\qquad Q=g-\tfrac{v^2}{2}c,
\end{gather*}
the weighted errors, the count-only margin, the rule contrast from \cref{eq:rule}, and a gate--count balance. Let $P_b$ be the
probability that the prediction with match bit $b$ assigns to its target.

\begin{lemma}[Derivatives]
\label{lem:J-geometry}
\begin{enumerate}[label=(\alph*),leftmargin=*,itemsep=2pt,topsep=3pt]
\item $\pd{\mathcal L}{c}=\mathrm{err}_1-\mathrm{err}_0$, $\quad\pd{\mathcal L}{g}=r_1\mathrm{err}_1-r_0\mathrm{err}_0=\tfrac{1}{2}\bigl(\pd{\mathcal L}{c}+\tau\,\mathrm{err}\bigr)$, $\quad\pd{\mathcal L}{a}=g\,\dd{\sigma(a)}{a}\,\mathrm{err}$.
\item $0<\phi<1$; $\phi(m;s)\ge\frac{1}{2}\sigma(-m)$; $\phi(m;s)\ge\frac{1}{2}(1-P_b)$ at $m=m_b$; $m\mapsto\phi(m;s)$ is nonincreasing; and
$\phi(m-2x;s)\le e^{2x}\phi(m;s)$ for $x\ge0$.
\item $F(m;s)\ge\ell(m)\ge-m$.
\item $\ddt{\mathcal L}=-\rho\initscale^2(\pd{\mathcal L}{c})^2-\rho\initscale^2v^2(\pd{\mathcal L}{g})^2-\frac{\initscale^2}{2N}(\pd{\mathcal L}{s})^2-\frac{\initscale^4S}{d_{\mathrm h}}(\pd{\mathcal L}{a})^2\le0$ (no $s$ term if $K=2$).
\end{enumerate}
\end{lemma}

\begin{proof}
(a) Differentiate the loss and margins in \cref{lem:J-reduce}, using
$\pd{F}{m}=-\phi$ and $r_0+r_1=1$.

(b) With $A=e^{-m}$ and $B=Ne^{-s-m/2}$, $\phi=\frac{A+B/2}{1+A+B}$: the probability of the other count plus half the background
probability. Hence $0<\phi<1$ and
 $\phi\ge\frac{1}{2}\cdot\frac{A+B}{1+A+B}=\frac{1}{2}(1-P_b)\ge\frac{1}{2}\cdot\frac{A}{1+A}=\frac{1}{2}\sigma(-m).$ 
Convexity of log-sum-exp \citep[Section~3.1.5]{boyd2004convex}, composed with these affine logits, makes $F(\cdot;s)$ convex and $\phi$ nonincreasing. Finally
 $\frac{d}{dm}\log\phi=-\frac{A+B/4}{A+B/2}+\frac{A+B/2}{1+A+B}>-1,$ 
and integration over $[m-2x,m]$ gives the last claim.

(c) Drop $Ne^{-s-m/2}$; $\log(1+e^{-m})\ge-m$. (d) is the chain rule with \cref{eq:J-flow}.
\end{proof}

\begin{lemma}[Loss levels]
\label{lem:J-levels}
Let $\mathcal L^\star=\frac{1}{2}p_{\min}\log2$.
\begin{enumerate}[label=(\alph*),leftmargin=*,itemsep=2pt,topsep=3pt]
\item If $\mathcal L<p_{\min}\log2$, then $P_0,P_1>\frac{1}{2}$: every prediction of $s_t$, $t\ge2$, has positive margin against all $K-1$ other counts.
\item If $P_0,P_1\ge\frac{1}{2}$, then $\mathrm{err}_0^2+\mathrm{err}_1^2\ge \frac{\mathcal L^2}{32}$.
\item If $\mathcal L>\mathcal L^\star$, then $\mathrm{err}_0^2+\mathrm{err}_1^2\ge \frac{p_{\min}^2}{267}$.
\end{enumerate}
\end{lemma}

\begin{proof}
(a) $\pi_b\ge p_{\min}$, so $-\log P_b=F(m_b;s)<\log2$.

(b) For $P\ge\frac{1}{2}$, $-\log P\le\frac{1-P}{P}\le2(1-P)$. By \cref{lem:J-geometry}(b), $F(m_b;s)\le4\phi(m_b;s)$, so $\mathcal L\le4\,\mathrm{err}$ and
$\mathrm{err}_0^2+\mathrm{err}_1^2\ge\frac{1}{2}\mathrm{err}^2\ge\frac{1}{32}\mathcal L^2$.

(c) If $P_b<\frac{1}{2}$, then $\mathrm{err}_b\ge\frac{\pi_b}{4}\ge\frac{p_{\min}}{4}$ by \cref{lem:J-geometry}(b); otherwise (b) applies. So
 $\mathrm{err}_0^2+\mathrm{err}_1^2\ge\min\Bigl\{\frac{p_{\min}^2}{16},\ \frac{(\mathcal L^\star)^2}{32}\Bigr\}=\frac{(\log2)^2}{128}\,p_{\min}^2\ge\frac{p_{\min}^2}{267}. \qedhere$ 
\end{proof}

\begin{lemma}[Gradient lower bound]
\label{lem:J-gradient}
$(\pd{\mathcal L}{c})^2+(\pd{\mathcal L}{g})^2\ge\frac{1}{3}\tau^2(\mathrm{err}_0^2+\mathrm{err}_1^2)$. If $v^2\le1$, then
 $\ddt{\mathcal L}\le-\tfrac{1}{3}\rho\initscale^2v^2\tau^2\,(\mathrm{err}_0^2+\mathrm{err}_1^2).$ 
\end{lemma}

\begin{proof}
By \cref{lem:J-geometry}(a), $(\pd{\mathcal L}{c},\pd{\mathcal L}{g})^\top=\mathbf C_{\mathrm{grad}}(\mathrm{err}_0,\mathrm{err}_1)^\top$ with $\mathbf C_{\mathrm{grad}}=\bigl(\begin{smallmatrix}-1&1\\-r_0&r_1\end{smallmatrix}\bigr)$. Its singular values $s_1\ge s_2$
satisfy
 \[
\begin{gathered}
s_1s_2=|\det\mathbf C_{\mathrm{grad}}|=|\tau|,\\ s_1\le\|\mathbf C_{\mathrm{grad}}\|_F=(2+r_0^2+r_1^2)^{1/2}\le\sqrt3,\\ s_2\ge\frac{|\tau|}{\sqrt3} .
\end{gathered}
\] 
The second claim follows from \cref{lem:J-geometry}(d).
\end{proof}

\paragraph{Reflection.} The map $(c,g,s,a)\mapsto(-c,-g,s,-a)$ exchanges $m_0$ and $m_1$, and so maps $\mathcal L$ and \cref{eq:J-flow} for $p$ to those
for $1-p$. On parameters it is $(u_0,u_1,w,q,k)\mapsto(u_1,u_0,-w,-q,k)$, i.e.\ $q(0)\mapsto-q(0)$. It preserves $q(0)k(0)(1-2p)$, $\qkbalance$ and the
conclusion of \cref{thm:matched}. Hence we assume $p<\frac{1}{2}$, and set
 $h^\star=\log\frac{1-p}{p}>0,\qquad \mu=\frac{p(1-2p)}{14}.$ 
The initialization is \emph{aligned} if $q(0)k(0)<0$, i.e.\ $a(0)<0$.

\begin{lemma}[Drift identities]
\label{lem:J-drift}
\begin{enumerate}[label=(\alph*),leftmargin=*,itemsep=2pt,topsep=3pt]
\item $\operatorname{sign}\ddt{a}=-\operatorname{sign}g$.
\item $\ddt{Q}=-\frac{1}{2}\rho\initscale^2v^2\,\tau\,\mathrm{err}$, so $Q$ is nondecreasing while $a\le0$.
\item $(1+\frac{v^2}{4})\,g=Q+\frac{v^2}{2}\,h$.
\item $\ddt{h}=-\rho\initscale^2(1+\frac{v^2}{4})\pd{\mathcal L}{c}-\frac{1}{4}\rho\initscale^2v^2\,\tau\,\mathrm{err}$.
\item If $h\le\frac{h^\star}{2}$, then $\pd{\mathcal L}{c}\le0$. If $h\le\frac{h^\star}{4}$ and $|R|\le1$, then $\pd{\mathcal L}{c}\le-\mu$.
\end{enumerate}
\end{lemma}

\begin{proof}
(a) By \cref{eq:J-flow} and \cref{lem:J-geometry}(a), $\ddt{a}=-\frac{\initscale^4S}{d_{\mathrm h}}\,g\,\dd{\sigma(a)}{a}\,\mathrm{err}$ with $\dd{\sigma(a)}{a},\mathrm{err}>0$ and $S=q^2+k^2>0$ (if $\qkbalance=0$, then $S=\frac{2\sqrt{d_{\mathrm h}}|a|}{\initscale^2}$ and $a$ never reaches $0$ by \cref{rem:J-balance}).

(b), (d) By \cref{eq:J-flow} and \cref{lem:J-geometry}(a), $\ddt{Q}=-\rho\initscale^2v^2\bigl(\pd{\mathcal L}{g}-\frac{1}{2}\pd{\mathcal L}{c}\bigr)$ and
$\ddt{h}=\ddt{c}+\frac{1}{2}\ddt{g}$. (c) Substitute $c=h-\frac{g}{2}$.

(e) $m_1=m_0-2h$ and $\pd{\mathcal L}{c}=p\,\phi(m_1;s)-(1-p)\,\phi(m_0;s)$. By \cref{lem:J-geometry}(b),
\begin{alignat*}{2}
 \phi(m_1;s)&\le e^{2\max(h,0)}\phi(m_0;s) &\qquad& \why{$\phi$ nonincreasing; $\phi(m-2x)\le e^{2x}\phi(m)$}\\
 \pd{\mathcal L}{c}&\le-\phi(m_0;s)\bigl((1-p)-p\,e^{2\max(h,0)}\bigr),
\end{alignat*}
and $p\,e^{h^\star}=1-p$ gives the first claim. If $h\le\frac{h^\star}{4}$, then $p\,e^{2\max(h,0)}\le p\,e^{h^\star/2}=\sqrt{p(1-p)}$ and
\begin{alignat*}{2}
 (1-p)-\sqrt{p(1-p)}&=\frac{\sqrt{1-p}\,(1-2p)}{\sqrt{1-p}+\sqrt p}\ge\frac{1-2p}{2} &\qquad& \why{$\sqrt{1-p}\ge\tfrac{1}{\sqrt2}\ge\tfrac{1}{2}(\sqrt{1-p}+\sqrt p)$}\\
 \phi(m_0;s)&\ge\tfrac{1}{2}\sigma\bigl(-\tfrac{h^\star}{4}-\tfrac{1}{2}\bigr) && \why{$m_0=h+\tfrac{R}{2}\le\tfrac{h^\star}{4}+\tfrac{1}{2}$}\\
 &\ge\tfrac{1}{4}e^{-1/2}\Bigl(\frac{p}{1-p}\Bigr)^{1/4}\ge\frac{p}{7} && \why{$\sigma(-z)\ge\tfrac{1}{2}e^{-z}$}
\end{alignat*}
using $p^{1/4}\ge p$ and $\frac{1}{4}e^{-1/2}\ge\frac{1}{7}$. Multiply.
\end{proof}

\begin{lemma}[Alignment is kept]
\label{lem:J-align}
Let $t_0\ge0$ with
 $a(t_0)<0,\qquad g(t_0)\ge0,\qquad \bigl(1+\tfrac{v^2}{4}\bigr)g(t_0)\ge\tfrac{v^2}{2}\bigl(h(t_0)-\tfrac{h^\star}{2}\bigr)_+ .$ 
Then for all $t\ge t_0$: $g(t)\ge0$, $a(t)\le a(t_0)$, and $h(t)\ge h_0:=\min\{h(t_0),\frac{h^\star}{2}\}$.
\end{lemma}

\begin{proof}
Let $t'=\sup\{t\ge t_0:a<0\text{ on }[t_0,t]\}>t_0$. On $[t_0,t')$, $\tau<0$ and $\mathrm{err}>0$.

\emph{$h\ge h_0$.} Where $h\le\frac{h^\star}{2}$, $\ddt{h}>0$ by \cref{lem:J-drift}(d), (e). If $h(t_1)<h_0$ for some $t_1\in(t_0,t')$, let
$t_2=\sup\{t\le t_1:h(t)\ge h_0\}$. Then $h<h_0\le\frac{h^\star}{2}$ on $(t_2,t_1]$, so $h(t_1)>h(t_2)=h_0$, a contradiction.

\emph{$g\ge0$.} By \cref{lem:J-drift}(b), (c),
\begin{alignat*}{2}
 \bigl(1+\tfrac{v^2}{4}\bigr)g(t)&=\bigl(1+\tfrac{v^2}{4}\bigr)g(t_0)+\bigl(Q(t)-Q(t_0)\bigr) &\qquad&\\
 &\qquad+\tfrac{v^2}{2}\bigl(h(t)-h(t_0)\bigr)\\
 &\ge\bigl(1+\tfrac{v^2}{4}\bigr)g(t_0)-\tfrac{v^2}{2}\bigl(h(t_0)-h_0\bigr) && \why{$Q\uparrow$, $h\ge h_0$}\\
 &\ge0 && \why{$h(t_0)-h_0=(h(t_0)-\tfrac{h^\star}{2})_+$}
\end{alignat*}

\emph{$a\le a(t_0)$.} By \cref{lem:J-drift}(a), $\ddt{a}\le0$ on $[t_0,t')$. If $t'<\infty$, then $a(t')\le a(t_0)<0$ and, by continuity, $a<0$ on a
neighbourhood of $t'$, contradicting the definition of $t'$. So $t'=\infty$.
\end{proof}

\subsection{Conditions and scales}

\paragraph{Aligned threshold and scales.}\label{app:J-constants}
After reflection assume $p<1/2$. We use only $0<\varepsilon\le1/2$ and $q(0)k(0)<0$.
Let $a_{\mathrm{sc}}=\varepsilon^4/\sqrt{d_{\mathrm h}}$. The seed implies
\begin{equation}\label{eq:J-scales}
\begin{gathered}\tfrac14\varepsilon^2\le v^2\le4\varepsilon^2\le1,\quad
\tfrac14 a_{\mathrm{sc}}\le|a(0)|\le4a_{\mathrm{sc}}\le\tfrac14,\\
\sigma'(a)\ge\tfrac16\ (|a|\le1),\quad
\tanh(x/2)\ge x/4\ (0\le x\le2).\end{gathered}
\end{equation}
\subsection{Aligned gate growth and persistent margins}

The next two lemmas give a persistent readout floor and a time bound for $|a|$ to reach $1$.

Let $T_h=\frac{h^\star}{4\rho\initscale^2\mu}$ and $g_{\mathrm{low}}=\min\{1,\frac{v^2h^\star}{20}\}\ge\frac{\initscale^2}{80}\min\{1,h^\star\}$.

\begin{lemma}[Gate floor]
\label{lem:J-floor}
Let $t_0$ satisfy the hypotheses of \cref{lem:J-align}, with moreover $h(t_0)\ge0$ and $Q(t_0)\ge-\frac{v^2h^\star}{16}$. Then $g(t)\ge g_{\mathrm{low}}$ for all
$t\ge t_0+T_h$.
\end{lemma}

\begin{proof}
On $[t_0,\infty)$, $a<0$, $Q\ge Q(t_0)$ and $h\ge0$ (\cref{lem:J-align}, \cref{lem:J-drift}(b)). At any $s\ge t_0$ with $g(s)<g_{\mathrm{low}}$,
\begin{alignat*}{2}
 \tfrac{v^2}{2}h(s)&=\bigl(1+\tfrac{v^2}{4}\bigr)g(s)-Q(s)<\tfrac{5}{4}\cdot\tfrac{v^2h^\star}{20}+\tfrac{v^2h^\star}{16}=\tfrac{v^2h^\star}{8} &\qquad& \why{\cref{lem:J-drift}(c), $v^2\le1$}\\
 |R(s)|&=g(s)|\tau(s)|<g_{\mathrm{low}}\le1,
\end{alignat*}
so $h(s)<\frac{h^\star}{4}$, and by \cref{lem:J-drift}(d), (e) and $\tau(s)<0$,
 $\ddt{h}(s)\ge\rho\initscale^2\mu,\qquad \frac{d}{dt}\Bigl[\bigl(1+\tfrac{v^2}{4}\bigr)g\Bigr](s)=\ddt{Q}(s)+\tfrac{v^2}{2}\ddt{h}(s)>0 .$ 
Suppose $g(t)<g_{\mathrm{low}}$ for some $t\ge t_0+T_h$, and let $s_0=\sup\{s\in[t_0,t]:g(s)\ge g_{\mathrm{low}}\}$, or $s_0=t_0$ if there is none. On $(s_0,t]$, $g$
increases, so $g(t)>g(s_0)$; if $s_0>t_0$ then $g(s_0)=g_{\mathrm{low}}$, a contradiction. So $s_0=t_0$, and on $(t_0,t]$, $h$ increases at rate at least
$\rho\initscale^2\mu$: $h(t)\ge\rho\initscale^2\mu\,T_h=\frac{h^\star}{4}$, contradicting $h(t)<\frac{h^\star}{4}$.
\end{proof}

\begin{lemma}[Gate growth]
\label{lem:J-grow}
Let $t_0\ge0$ with $a(t_0)\le-\frac{a_{\mathrm{sc}}}{8}$, and suppose $g\ge g_{\mathrm{low}}$ on $[t_0,\infty)$. Then $t_3=\inf\{t\ge t_0:\mathcal L(t)\le\mathcal L^\star\}$ satisfies
 $t_3-t_0\le\frac{4080\sqrt{d_{\mathrm h}}\,\log\frac{8\sqrt{d_{\mathrm h}}}{\initscale^4}}{p_{\min}\min\{1,h^\star\}\,\initscale^4}+\frac{8010\,K\log K}{p_{\min}^2\initscale^4} .$ 
\end{lemma}

\begin{proof}
Since $g\ge0$, $a$ is nonincreasing (\cref{lem:J-drift}(a)), so $|a|$ is nondecreasing on $[t_0,\infty)$. Let $t_1=\inf\{t\ge t_0:|a(t)|\ge1\}$. On
$[t_0,\min\{t_1,t_3\})$,
\begin{alignat*}{2}
 \mathrm{err}&\ge\bigl(\mathrm{err}_0^2+\mathrm{err}_1^2\bigr)^{1/2}\ge\frac{p_{\min}}{17} &\qquad& \why{\cref{lem:J-levels}(c)}\\
 \frac{d|a|}{dt}&=\frac{\initscale^4S}{d_{\mathrm h}}\,g\,\dd{\sigma(a)}{a}\,\mathrm{err}\ge\frac{2\initscale^2|a|}{\sqrt{d_{\mathrm h}}}\cdot\frac{\initscale^2\min\{1,h^\star\}}{80}\cdot\frac{1}{6}\cdot\frac{p_{\min}}{17} && \why{$S\ge\frac{2\sqrt{d_{\mathrm h}}|a|}{\initscale^2}$, $|a|\le1$}\\
 &=\frac{\initscale^4\min\{1,h^\star\}\,p_{\min}}{4080\sqrt{d_{\mathrm h}}}\,|a|,
\end{alignat*}
so $\min\{t_1,t_3\}-t_0\le\frac{4080\sqrt{d_{\mathrm h}}}{\initscale^4\min\{1,h^\star\}p_{\min}}\log\frac{8}{a_{\mathrm{sc}}}$. If $t_1<t_3$, then on $[t_1,t_3)$, $|a|\ge1$ and
\begin{alignat*}{2}
 \ddt{\mathcal L}&\le-\tfrac{1}{3}\rho\initscale^2v^2\tanh^2(\tfrac{1}{2})\,(\mathrm{err}_0^2+\mathrm{err}_1^2) &\qquad& \why{\cref{lem:J-gradient}}\\
 &\le-\frac{\initscale^4p_{\min}^2}{8010\,K} && \why{\cref{lem:J-levels}(c), $\tanh^2\tfrac{1}{2}\ge\tfrac{1}{5}$}
\end{alignat*}
and $t_3-t_1\le\frac{8010K\log K}{p_{\min}^2\initscale^4}$, since $\mathcal L(t_1)\le\log K$.
\end{proof}

\begin{lemma}[Descent]
\label{lem:J-descent}
Suppose $|a(t)|\ge\frac{a_{\mathrm{sc}}}{8}$ for all $t\ge t_0$, and let $t_3=\inf\{t\ge t_0:\mathcal L(t)\le \mathcal L^\star\}$. Then
\[
 t_3-t_0\le\frac{1.7\cdot10^6\,K\log K\,d_{\mathrm h}}{p_{\min}^2\initscale^{12}},
\]
every prediction of $s_t$, $t\ge2$, has positive margin against all $K-1$ other counts for all $t\ge t_3$, and
$\mathcal L(t)\le\frac{2\cdot10^5Kd_{\mathrm h}}{\initscale^{12}(t-t_3)}$ for $t>t_3$.
\end{lemma}

\begin{proof}
For $t\ge t_0$,
\begin{alignat*}{2}
 \tau^2&\ge\tanh^2(\tfrac{a_{\mathrm{sc}}}{16})\ge\frac{a_{\mathrm{sc}}^2}{1024} &\qquad& \why{\cref{eq:J-scales}}\\
 \ddt{\mathcal L}&\le-\tfrac{1}{3}\rho\initscale^2v^2\tau^2\,(\mathrm{err}_0^2+\mathrm{err}_1^2) && \why{\cref{lem:J-gradient}}\\
 &\le-\frac{\initscale^{12}}{6144\,K\,d_{\mathrm h}}\,(\mathrm{err}_0^2+\mathrm{err}_1^2) && \why{$\rho=\tfrac{2}{K}$, $v^2\ge\tfrac{1}{4}\initscale^2$}
\end{alignat*}
While $\mathcal L>\mathcal L^\star$,
\begin{alignat*}{2}
 \ddt{\mathcal L}&\le-\frac{\initscale^{12}p_{\min}^2}{1640448\,K\,d_{\mathrm h}} &\qquad& \why{\cref{lem:J-levels}(c)}\\
 t_3-t_0&\le\frac{1640448\,K\log K\,d_{\mathrm h}}{p_{\min}^2\initscale^{12}} && \why{$\mathcal L(t_0)\le\mathcal L(0)=\log K$}
\end{alignat*}
For $t\ge t_3$, $\mathcal L<p_{\min}\log2$, so all margins are positive and $P_0,P_1>\frac{1}{2}$ (\cref{lem:J-levels}(a)), and
\begin{alignat*}{2}
 \ddt{\mathcal L}&\le-\frac{\initscale^{12}}{196608\,K\,d_{\mathrm h}}\,\mathcal L^2 &\qquad& \why{\cref{lem:J-levels}(b)}\\
 \frac{d}{dt}\frac{1}{\mathcal L}&\ge\frac{\initscale^{12}}{196608\,K\,d_{\mathrm h}} && \qedhere
\end{alignat*}
\end{proof}

\begin{proof}[Process part of \cref{thm:matched}]
\emph{Reduction.} By the reflection, $p<\frac{1}{2}$; then $p_{\min}=p$.

\emph{Aligned start.} $a(0)\le-\frac{a_{\mathrm{sc}}}{4}<0$ and $g(0)=h(0)=0$, so \cref{lem:J-align} applies at $t_0=0$: $a(t)\le-\frac{a_{\mathrm{sc}}}{4}$ for all $t$.
\Cref{lem:J-descent} with $t_0=0$ gives positive margins for $t\ge t_3$ and $\Lstep\to0$. \Cref{lem:J-floor} applies at $t_0=0$ ($h(0)=Q(0)=0$), and
\cref{lem:J-grow} with $t_0=T_h$ gives
 $t_3\le T_h+\frac{4080\sqrt{d_{\mathrm h}}\,\log\frac{8\sqrt{d_{\mathrm h}}}{\initscale^4}}{p\min\{1,h^\star\}\,\initscale^4}+\frac{8010\,K\log K}{p^2\initscale^4} .$ 


The explicit bound in \cref{eq:focused-time} is precisely this aligned estimate.
The loss is nonincreasing after crossing $\mathcal L^\star$; the strict threshold
in \cref{lem:J-levels}(a) therefore makes the margins persistent. The decay
estimate of \cref{lem:J-descent} proves $\Lstep\to0$.
\end{proof}
\paragraph{Scalar balance.}\label{rem:J-balance}
The identity $q^2-k^2=\Delta_{qk}$ is invariant. Even when it vanishes,
alignment persists: the equation for $a$ has the form $\dot a=-f(t)a$
on a finite time interval, with continuous $f$. Thus a nonzero aligned
$a(0)$ never reaches zero in finite time.
\paragraph{Readout restriction.}\label{rem:J-line}
The line $\mathbb R(I-R)$ is part of the trained family, not a conclusion
about unrestricted gate readouts. For $K=2$ it removes only a common-logit
component. For $K\ge3$, unrestricted circulant gate readouts introduce
additional background contributions to the gate derivative; this proof does
not cover them.
